% SPDX-License-Identifier: GPL-3.0-or-later
% Copyright (C) 2026 Alessandro Greco
% --- Adattamento delle figure alla colonna piu' stretta del 7x10 ------------
% Le figure TikZ sono tarate per il testo A4 (piu' largo); nel 7x10 alcune
% sforano il margine. Le riduciamo automaticamente a \linewidth (max width:
% rimpicciolisce solo quelle troppo larghe, lascia intatte le altre). Vale solo
% per questa edizione: i file dei capitoli e l'edizione A4 non cambiano.
\usepackage{etoolbox}
\usepackage{adjustbox}
% Only print builds set this switch; A4 and EPUB retain their original diagrams.
\def\PrintEdition{1}
\BeforeBeginEnvironment{tikzpicture}{\begin{adjustbox}{max width=0.96\linewidth}}
\AfterEndEnvironment{tikzpicture}{\end{adjustbox}}
\lstset{numbersep=6pt,xleftmargin=15pt,xrightmargin=16pt,breakatwhitespace=false,frame=none,framerule=0pt}  % niente bordo (sfondo rientrato): evita il filetto che sporgeva

% --- Testatina adatta alla colonna stretta ---------------------------------
% I titoli di capitolo lunghi, in maiuscolo, sforavano la testatina nel 7x10.
% Qui la rendiamo piu' piccola e senza maiuscolo (solo questa edizione).
\usepackage{fancyhdr}
\microtypesetup{protrusion=false}
\hypersetup{hidelinks}  % annotations are disabled by NoHyper below
\pagestyle{fancy}
\AtBeginDocument{\setlength{\headwidth}{\textwidth}}  % dopo geometry: testatina = larghezza del testo
\fancyhf{}
\fancyhead[LE,RO]{\footnotesize\thepage}
\fancyhead[RE]{\adjustbox{max width=0.88\textwidth}{\footnotesize\nouppercase{\itshape\leftmark}}}
\fancyhead[LO]{\adjustbox{max width=0.88\textwidth}{\footnotesize\nouppercase{\itshape\rightmark}}}
\renewcommand{\headrulewidth}{0pt}
\fancypagestyle{plain}{\fancyhf{}\fancyfoot[C]{\footnotesize\thepage}\renewcommand{\headrulewidth}{0pt}}
% \pagestyle{fancy} ridefinisce \chaptermark senza distinguere \frontmatter: la seconda
% pagina della Prefazione portava "Capitolo 0. Prefazione". Qui si ripristina il test
% \if@mainmatter della classe book: numero solo nel corpo del libro, \@chapapp
% (Capitolo/Appendice) come prima.
\makeatletter
\renewcommand{\chaptermark}[1]{\markboth{\MakeUppercase{\ifnum\c@secnumdepth>\m@ne\if@mainmatter\@chapapp\ \thechapter. \ \fi\fi #1}}{}}
\makeatother

% --- Formato pagina 7 x 10 in ----------------------------------------------
% 7x10 no-bleed, 151--400 pages: 1in inside, 0.53in outside (0.03in reserve for italic ink overhang).
\geometry{
  paperwidth=7in, paperheight=10in,
  inner=1in, outer=0.53in,
  top=0.75in, bottom=0.75in,
  headheight=14pt, headsep=6pt, footskip=16pt,
  marginparwidth=0pt, marginparsep=0pt,
}
% Un filo di elasticita' in piu' per la giustificazione sulla colonna stretta,
% cosi' l'algoritmo preferisce allargare gli spazi anziche' sforare il margine.
\setlength{\emergencystretch}{4em}
% Spaziatura fra colonne leggermente ridotta (default 6pt): sulla colonna piu'
% stretta del 7x10 evita che le tabelle a larghezza fissa (tarate per l'A4)
% sforino il margine esterno. Vale solo per questa edizione.
\setlength{\tabcolsep}{4pt}

% Per la stampa i collegamenti restano neri e senza bordo: hidelinks e' gia'
% impostato in preambolo.tex, quindi non serve altro.

% Le pagine bianche inserite da \cleardoublepage (verso del frontespizio, verso
% prima di un capitolo con openright) devono essere del tutto vuote, senza numero
% ne' testatina: piu' pulito per un libro stampato.
\makeatletter
\def\cleardoublepage{\clearpage\if@twoside\ifodd\c@page\else
  \hbox{}\thispagestyle{empty}\newpage\if@twocolumn\hbox{}\newpage\fi\fi\fi}
\makeatother


% Keep hyperref metadata but omit clickable annotations on paper.
\AtBeginDocument{\begin{NoHyper}}
\AtEndDocument{\end{NoHyper}}
% Missing characters are errors, not silently accepted omissions.
\tracinglostchars=3
% Run-in headings must finish before a listing, even without its frame.
\makeatletter
\AddToHook{env/lstlisting/before}{\if@noskipsec\leavevmode\par\fi}
\makeatother
% Fixed-width tables can only shrink; multipage tables remain breakable.
\BeforeBeginEnvironment{tabular}{\begin{adjustbox}{max width=\linewidth}}
\AfterEndEnvironment{tabular}{\end{adjustbox}}
% The glossary is the only longtable in each edition. Adapt column widths,
% preserving every cell and repeating header; never box a multipage table.
\let\PrintLongtable\longtable
\providecommand{\PrintGlossaryColumns}{@{}>{\raggedright\arraybackslash}p{0.29\linewidth}>{\raggedright\arraybackslash}p{0.57\linewidth}r@{}}
\RenewDocumentCommand\longtable{m}{%
  \expandafter\PrintLongtable\expandafter{\PrintGlossaryColumns}}


% hyperxmp keeps the author in XMP; retain the standard PDF Info entry too.
\AtBeginDocument{\pdfextension info{/Author (Alessandro Greco)}}

% Break long inline concatenations at their existing mathematical separator.
% No \left\| or \right\| delimiters occur in the unchanged book sources.
\let\PrintConcat\|
\renewcommand{\|}{\allowbreak\PrintConcat}
\usepackage{xurl}
% Single-line unnumbered displays in these sources can exceed the print measure.
% Preserve every symbol; shrink only the displays that are actually too wide.
\let\PrintDisplayBegin\[
\let\PrintDisplayEnd\]
\long\def\[#1\]{\PrintDisplayBegin
  \adjustbox{max width=\linewidth}{$\displaystyle #1$}%
  \PrintDisplayEnd}

% Additional break opportunities in file names/API calls, without printed hyphens.
% The A4 \file definition remains untouched.
\ExplSyntaxOn
\RenewDocumentCommand{\file}{m}
  {
    \group_begin:
    \tl_set:Nn \l__libro_file_tl {#1}
    \tl_replace_all:Nnn \l__libro_file_tl { / } { / \allowbreak }
    \tl_replace_all:Nnn \l__libro_file_tl { . } { . \allowbreak }
    \tl_replace_all:Nnn \l__libro_file_tl { :: } { :: \allowbreak }
    \tl_replace_all:Nnn \l__libro_file_tl { ( } { ( \allowbreak }
    \cs_set_eq:NN \__libro_trattino_basso: \_
    \cs_set:Npn \_ { \__libro_trattino_basso: \allowbreak }
    \texttt { \tl_use:N \l__libro_file_tl }
    \group_end:
  }
\ExplSyntaxOff

% More line-breaking flexibility for the narrower measure.
\tolerance=2000
\indexsetup{othercode=\small}
